\section*{Capítulo \ref{ch:g10:biodiversity-scales} --- A biodiversidade em todas as escalas}

\begin{solution}{exo:g10:biodiversity-scales:1}
Os \omterm{def:g10:biodiversity-scales:biodiversity}{ecossistemas} (uma lagoa, uma floresta de faias); as \omterm{def:g10:biodiversity-scales:species}{espécies} (as
rãs, os tritões e os juncos da lagoa); os \omterm{def:g10:universal-dna:gene}{genes} (os \omterm{def:g10:universal-dna:gene}{alelos} do \omterm{def:g10:universal-dna:gene}{gene} ABO
numa população humana).
\end{solution}

\begin{solution}{exo:g10:biodiversity-scales:2}
Um grupo de organismos parecidos que se reproduzem juntos e dão
descendentes férteis. O leão e o tigre dão híbridos estéreis, logo são
duas \omterm{def:g10:biodiversity-scales:species}{espécies}.
\end{solution}

\begin{solution}{exo:g10:biodiversity-scales:3}
$\num{1000000}/6400 \approx 156$: cerca de 150 vezes mais.
\end{solution}

\begin{solution}{exo:g10:biodiversity-scales:4}
A fração de todas as cópias de um \omterm{def:g10:universal-dna:gene}{gene} numa população que são um dado
\omterm{def:g10:universal-dna:gene}{alelo}. B $= 1 - 0.66 - 0.27 = 0.07$.
\end{solution}

\begin{solution}{exo:g10:biodiversity-scales:5}
Cinco; a maior foi a extinção do fim do Permiano, há 252 milhões de
anos, que eliminou cerca de 95\% das \omterm{def:g10:biodiversity-scales:species}{espécies} marinhas.
\end{solution}

\begin{solution}{exo:g10:biodiversity-scales:6}
O segundo: riqueza igual, mas uma abundância relativa bem mais
uniforme. A riqueza sozinha não os distingue; as abundâncias relativas
sim.
\end{solution}

\begin{solution}{exo:g10:biodiversity-scales:7}
O campo é o mesmo; a colega amostrou dez vezes mais. As \omterm{def:g10:biodiversity-scales:species}{espécies}
encontradas aumentam com o esforço até a curva se achatar, de modo que
contagens feitas com esforços diferentes não são comparáveis --- os 8
do aluno são uma subamostragem, não um campo mais pobre.
\end{solution}

\begin{solution}{exo:g10:biodiversity-scales:8}
Os poucos sobreviventes carregavam apenas uma fração dos \omterm{def:g10:universal-dna:gene}{alelos} da
\omterm{def:g10:biodiversity-scales:species}{espécie}; os descendentes deles, por mais numerosos que sejam,
compartilham esse conjunto reduzido. A \omterm{def:g10:biodiversity-scales:biodiversity}{diversidade genética} se perde no
gargalo e não é restaurada pelo número: a \omterm{def:g10:biodiversity-scales:species}{espécie} fica sem nada de
reserva contra uma doença nova ou um clima alterado (o caso do
guepardo).
\end{solution}

\begin{solution}{exo:g10:biodiversity-scales:9}
Cerca de 560 famílias antes, 300 depois: $260/560 \approx 45\%$ das
famílias perdidas. Uma família só desaparece quando todas as \omterm{def:g10:biodiversity-scales:species}{espécies}
dela desaparecem, de modo que muitas famílias sobreviveram por uma ou
duas \omterm{def:g10:biodiversity-scales:species}{espécies} e perderam o resto; a fração de \omterm{def:g10:biodiversity-scales:species}{espécies} perdidas (cerca
de 95\%) é bem maior que a fração de famílias.
\end{solution}

\begin{solution}{exo:g10:biodiversity-scales:10}
Toda população carrega os três \omterm{def:g10:universal-dna:gene}{alelos}, e os membros dela diferem entre
si em quais carregam; as diferenças entre populações são apenas
deslocamentos de frequência. A maior parte da \omterm{def:g10:biodiversity-scales:biodiversity}{diversidade genética} da
\omterm{def:g10:biodiversity-scales:species}{espécie} existe, portanto, dentro de cada população: dois vizinhos
diferem tipicamente em tantas posições quanto duas pessoas de lugares
distantes.
\end{solution}

\begin{solution}{exo:g10:biodiversity-scales:11}
\omterm{def:g10:biodiversity-scales:biodiversity}{Ecossistema}: a própria lagoa, com a comunidade dela. \omterm{def:g10:biodiversity-scales:species}{Espécies}: o tritão
ou a libélula que localmente só viviam ali. \omterm{def:g10:universal-dna:gene}{Genes}: os \omterm{def:g10:universal-dna:gene}{alelos} carregados
pela população de rãs da lagoa, que podiam diferir dos de outras
lagoas.
\end{solution}

\begin{solution}{exo:g10:biodiversity-scales:12}
Com uma duração de vida de cerca de $\num{2e6}$ anos, $\num{2e6}$
\omterm{def:g10:biodiversity-scales:species}{espécies} dão cerca de uma extinção por ano. Mil vezes mais são cerca de
1000 \omterm{def:g10:biodiversity-scales:species}{espécies} por ano, $\num{100000}$ por século --- 5\% das \omterm{def:g10:biodiversity-scales:species}{espécies}
conhecidas em cem anos.
\end{solution}

\begin{solution}{exo:g10:biodiversity-scales:13}
A extinção esvaziou modos de vida --- grandes herbívoros, grandes
predadores, voadores, nadadores --- que os dinossauros tinham ocupado.
As linhagens de mamíferos sobreviventes, livres desses concorrentes e
com a \omterm{def:g10:biodiversity-scales:biodiversity}{diversidade genética} própria a que recorrer, se adaptaram aos
papéis vagos e se dividiram em novas \omterm{def:g10:biodiversity-scales:species}{espécies}: a perda abriu as
oportunidades que a nova diversidade preencheu.
\end{solution}

\begin{solution}{exo:g10:biodiversity-scales:14}
O trigo cultivado perdeu a maior parte dos \omterm{def:g10:universal-dna:gene}{alelos} dos ancestrais dele.
Os parentes silvestres os conservam: \omterm{def:g10:universal-dna:gene}{alelos} de resistência a uma
doença, à seca, ao sal, que nenhuma variedade atual carrega. Cruzá-los
com o trigo devolve a opção; uma vez extinto um parente silvestre, os
\omterm{def:g10:universal-dna:gene}{alelos} dele se foram para sempre.
\end{solution}

\begin{solution}{exo:g10:biodiversity-scales:15}
A extinção é natural, mas a cerca de uma \omterm{def:g10:biodiversity-scales:species}{espécie} por ano; a taxa atual
é centenas a mil vezes mais rápida, comparável a uma \omterm{prop:g10:biodiversity-scales:changes}{extinção em massa},
e a recuperação daquelas levou milhões de anos --- mais tempo do que a
nossa \omterm{def:g10:biodiversity-scales:species}{espécie} existe. Enquanto isso, os \omterm{def:g10:biodiversity-scales:biodiversity}{ecossistemas} que estão sendo
perdidos fornecem alimento, polinização, solo, purificação da água e as
reservas genéticas das nossas lavouras. O precedente mostra que a perda
é real e que a recuperação não está disponível em nenhuma escala de
tempo humana.
\end{solution}

\begin{solution}{pb:g10:biodiversity-scales:1}
\textbf{1.} 31 e 9 \omterm{def:g10:biodiversity-scales:species}{espécies}: razão $31/9 \approx 3.4$.

\textbf{2.} A fazenda A: nenhuma \omterm{def:g10:biodiversity-scales:species}{espécie} acima de 22\%, contra uma a
84\% na fazenda B. A uniformidade diz como os indivíduos se dividem
entre as \omterm{def:g10:biodiversity-scales:species}{espécies} encontradas; uma comunidade rica mas desigual é
dominada por uma \omterm{def:g10:biodiversity-scales:species}{espécie} e o resto é raro.

\textbf{3.} Os ganhos diminuem (7, 4, 2): a curva está se achatando, de
modo que a amostra está perto de completa; mais algumas parcelas
acrescentariam no máximo uma ou duas \omterm{def:g10:biodiversity-scales:species}{espécies}.

\textbf{4.} A curva da fazenda B chegou a 9 na 15ª parcela e não
acrescentou nada depois: uma curva de amostragem completa se achata num
patamar, que a fazenda B alcançou e a fazenda A quase alcançou.

\textbf{5.} Muitas plantas só são visíveis ou identificáveis em flor, e
os insetos e as aves variam com a estação; um mês diferente mudaria as
contagens por razões sem relação com as cercas vivas.

\textbf{6.} Insetos $46/14 \approx 3.3$; aves $17/5 = 3.4$. Os três
grupos dão um fator de cerca de 3: a mesma história.

\textbf{7.} As cercas vivas fornecem as plantas (folhas, flores, madeira
morta, abrigo) de que os insetos se alimentam e onde se reproduzem;
menos \omterm{def:g10:biodiversity-scales:species}{espécies} de plantas significam menos \omterm{def:g10:biodiversity-scales:species}{espécies} de insetos; os
insetos e os frutos e locais de ninho da cerca viva alimentam e abrigam
as aves, de modo que a diversidade de aves cai por sua vez.

\textbf{8.} Não: a produtividade mede apenas a lavoura. Uma cerca viva
pode abrigar o campo do vento, segurar o solo, alojar os polinizadores e
os predadores das pragas da lavoura, e os efeitos dela aparecem em anos
ruins ou ao longo de décadas, e não no número de uma safra.

\textbf{9.} Uma ave que nidifica em cerca viva, como um melro ou uma
toutinegra (ou um caracol de cerca viva, um percevejo do espinheiro):
ela precisa dos arbustos para nidificar, se alimentar ou se abrigar, e
a fazenda B não os oferece mais.

\textbf{10.} Os \omterm{def:g10:biodiversity-scales:biodiversity}{ecossistemas}: a própria cerca viva, um habitat, foi
removida, e as \omterm{def:g10:biodiversity-scales:species}{espécies} e os \omterm{def:g10:universal-dna:gene}{genes} se foram com ela; as perdas nas
outras escalas são consequências.

\textbf{11.} Fazenda A: cópias de y $2 \times 100 + 200 = 400$, cópias
de b $200 + 2 \times 100 = 400$, em 800: frequência de y 0.5. Fazenda
B: y $8 + 32 = 40$, b $32 + 128 = 160$, em 200: frequência de y 0.2.

\textbf{12.} Fazenda A: os dois \omterm{def:g10:universal-dna:gene}{alelos} igualmente frequentes e 50\% dos
caracóis carregam os dois; fazenda B: 32\% carregam os dois e b domina.
A fazenda A é mais diversa.

\textbf{13.} No barranco descoberto os caracóis amarelos são comidos com
mais frequência; geração após geração, menos cópias de y foram
passadas adiante, e a frequência de y caiu de cerca de 0.5 para 0.2.

\textbf{14.} A vantagem se inverteria: os caracóis amarelos
sobreviveriam melhor na cerca viva sombreada e a frequência de y
voltaria a subir ao longo das gerações --- desde que o \omterm{def:g10:universal-dna:gene}{alelo} ainda
esteja presente.

\textbf{15.} A da fazenda B: numa população de algumas centenas, um ano
ruim ou algumas mortes por acaso podem eliminar toda cópia do \omterm{def:g10:universal-dna:gene}{alelo}
mais raro; entre dezenas de milhares, o \omterm{def:g10:universal-dna:gene}{alelo} é carregado por milhares
de caracóis e o acaso não consegue apagá-lo.

\textbf{16.} Fazenda A: $31 + 46 + 17 = 94$; fazenda B:
$9 + 14 + 5 = 28$; razão $94/28 \approx 3.4$.

\textbf{17.} Três por cento da terra abrigam o habitat de cerca de dois
terços das \omterm{def:g10:biodiversity-scales:species}{espécies} da fazenda.

\textbf{18.} Plantas $24/31 \approx 77\%$; insetos $35/46 \approx 76\%$;
aves $12/17 \approx 71\%$.

\textbf{19.} A \omterm{def:g10:biodiversity-scales:biodiversity}{diversidade genética}: as \omterm{def:g10:biodiversity-scales:species}{espécies} voltam de populações
vizinhas, mas os \omterm{def:g10:universal-dna:gene}{alelos} perdidos pelas populações locais (os caracóis da
fazenda B, por exemplo) só voltam se por acaso os imigrantes os
carregarem; um \omterm{def:g10:universal-dna:gene}{alelo} perdido não é recriado pelo replantio.

\textbf{20.} As cercas vivas multiplicaram a riqueza de \omterm{def:g10:biodiversity-scales:species}{espécies} da
fazenda por cerca de 3.4 em todos os grupos contados; uma cerca viva
replantada traz de volta o \omterm{def:g10:biodiversity-scales:biodiversity}{ecossistema} e, com o tempo, a maioria das
\omterm{def:g10:biodiversity-scales:species}{espécies}, mas não traz sozinha os \omterm{def:g10:universal-dna:gene}{alelos} que as populações locais
perderam.
\end{solution}
